\subsection{The exact central correction map and the fourth-layer local test}

\begin{statement}{A.205}{Proved}
{Compact correction fibres, geometric stabilizers, and the 240 cubed fourth-layer twists}
Fix an actual global section \(S\) of \(E_{r+1}\to G_\Q\),
where \(r\geq2\), and put \(s=pS\) in \(E_r\).
Write
\[
 Z=\widehat\Z(r+1)^{b_{r+1}},\qquad
 \mathcal H=H^1_{\rm cts}(G_\Q,Z).
\tag{A.205.1}
\]
The global geometric conjugacy classes of lifts above \([s]\)
are exactly the torsor \(S+\mathcal H\).
For each place \(v\), the actual local point and cusp sections
whose lower image is \([s_v]\) define the subset
\(\mathcal D_v\subset H^1(G_{\Q_v},Z)\) constructed below.
The full locally matching lift fibre is exactly
\[
 \{S+c:c\in\mathcal H,\quad
                   \operatorname{res}_v c\in\mathcal D_v
                      \text{ for every }v\}.
\tag{A.205.2}
\]
The group \(\mathcal H\) is compact.  The sets in this
intersection are closed.  Consequently this fibre is nonempty
exactly when every finite collection of its local conditions
has a solution in \(\mathcal H\).
The same criterion can be expressed through finitely many
places, coefficient primes and coefficient powers at a time.

Globally and at each finite place, the geometric centralizer
of the lower section is trivial.  At the real place its
possibly nontrivial centralizer and its exact translation
action are retained in (A.205.7).  The local allowed sets
retain all cusp labels: at a finite place they have at most
three elements, and they are singletons when the lower
section has an ordinary point image.

For the extension \(E_4\to E_3\) there is the exact
full-coefficient calculation
\[
 H^1_{\rm cts}(G_\Q,\widehat\Z(4))=\Z/240,
 \qquad \eta_4(g)=\frac{\chi(g)^4-1}{240}.
\tag{A.205.3}
\]
Thus every nonempty global lift fibre has exactly \(240^3\)
geometric section classes.  Restriction of its correction
group to every single finite place is injective.
The three coefficient projections give the actual isomorphism
\[
\begin{aligned}
 \Z/240&\xrightarrow{\ \sim\ }
 H^1(\Q_3,\Z_2(4))\oplus H^1(\Q_2,\Z_3(4))
                              \oplus H^1(\Q_2,\Z_5(4)),\\
 k&\longmapsto(k/3\bmod16,\ k/16\bmod3,\ k/16\bmod5).
\end{aligned}
\tag{A.205.4}
\]
For an all-place locally matching class-three section, A.204
supplies the initial \(S\).  After choosing its actual local
full point or cusp images at two and three, (A.205.4) selects
one global class-four lift matching precisely these three
coefficient projections, independently in the three Hall
coordinates.  Formula (A.205.2) is the remaining exact test
for full local matching.  The real comparison is automatic
for every such global class-four lift.

Failure of simultaneous matching has a witness using at most
\(240^3\) places.  If the lower section has an ordinary
image at one finite place, at most one global class-four
lift matches there; failure then has a witness using at
most that place and one further place.  No assertion that
the resulting finite test always succeeds is made.
\end{statement}

\begin{proof}
\emph{Actual lift torsors and the lower geometric centralizers.}
The kernel of \(E_{r+1}\to E_r\) is the central geometric
module \(Z\) in (A.205.1), with the actual cyclotomic
arithmetic action.  Every literal lift of \(s\) has the
unique form \(c(g)S(g)\), where \(c\) is a continuous
one-cocycle in \(Z\).  Conjugation by \(a\in Z\)
changes \(c\) by \(a-ga\), so quotienting by these
conjugators gives exactly \(H^1(G_\Q,Z)\).

For a global or local group \(G\), let
\[
 C_s(G)=\{u\in N_r:u\,s(g)\,u^{-1}=s(g)
                                  \text{ for every }g\in G\}.
\tag{A.205.5}
\]
At \(G=G_\Q\) or \(G_{\Q_p}\), this group is trivial.
Indeed a nonidentity element has a first nonzero Hall weight
\(j\leq r\), and at least one nonzero coefficient-prime
component in that weight.  Inner geometric conjugation
does not change that graded component.  Its invariance
would therefore give a nonzero invariant in
\(\Z_\ell(j)^{b_j}\).
For \(G_\Q\) the cyclotomic image is infinite.
For \(G_{\Q_p}\), when \(\ell\ne p\), arithmetic
Frobenius has cyclotomic value \(p\), with \(p^j-1\ne0\);
when \(\ell=p\), the local cyclotomic image is an infinite
open subgroup.  In each case some \(\chi_\ell(g)^j-1\)
is nonzero and cannot kill a nonzero element of a torsion-free
\(\Z_\ell\)-module.  This proves the claim.
It also proves that any geometric conjugacy between two
global lifts of the literal \(s\) is already central.
Aligning a lift above the class \([s]\) gives the same
global torsor, proving the first assertion.

The real stabilizer is not discarded.  Given
\(u\in C_s(G_\R)\), choose an actual geometric lift
\(\widetilde u\in N_{r+1}\), and define
\[
 d_u(g)=\widetilde uS(g)\widetilde u^{-1}S(g)^{-1}\in Z.
\tag{A.205.6}
\]
This is a continuous one-cocycle.  Replacing
\(\widetilde u\) by \(a\widetilde u\) changes it
by \(a-ga\).  Multiplying two lifted conjugators adds
their cocycle classes, since geometric conjugation on
\(Z\) is trivial.  We have therefore proved the actual
homomorphism and its action
\[
 \delta_s:C_s(G_\R)\longrightarrow H^1(G_\R,Z),
       \quad u\longmapsto[d_u],\qquad
             [c]\longmapsto[c]+\delta_s(u).
\tag{A.205.7}
\]
It is independent of the reference lift \(S\): a central
twist commutes with every geometric conjugator.  A continuous
Hall set lift of \(u\) proves continuity of this homomorphism.
The group \(C_s(G_\R)\) is closed in a compact geometric
group, so its image is compact.

\emph{The complete local allowed sets.}
For each actual full local point or cusp section \(\sigma_v\)
whose image in \(E_r\) is conjugate to \(s_v\), choose
a geometric conjugator making the lower representatives
equal and lift that conjugator to \(N_{r+1}\).
Denote the aligned class-\(r+1\) section by
\(\widetilde\sigma_v\).  The literal difference
\[
 c_{\sigma_v}(g)=\widetilde\sigma_v(g)S(g)^{-1}\in Z
\tag{A.205.8}
\]
is a continuous central cocycle.  Changes of the lifted
conjugator give a coboundary; changes of its lower alignment
translate its class by \(\delta_s(C_s(G_{\Q_v}))\),
with this subgroup zero at finite places.  Define
\(\mathcal D_v\) to be the union of these cosets for all
the actual local point and cusp images just specified.
The calculation proves both directions of (A.205.2).

These allowed sets are finite unions of closed cosets.
At a finite place the full multiplicative Kummer map on
\(\Q_p^\times\) is injective, by its valuation, root-of-unity
and principal-unit factors.  Thus the lower abelian pair
determines at most one ordinary centre.  For the cusp
parameters \(t,1-t,1/t\), the three literal abelian maps are
\[
 \tau\longmapsto(\tau,0),\qquad
 \tau\longmapsto(0,\tau),\qquad
 \tau\longmapsto(-\tau,\kappa(-1)-\tau).
\tag{A.205.9}
\]
Each fixes the completed tangent parameter uniquely when
the label is specified.  The cusp arithmetic decomposition
extension splits by that parameter, and its sections modulo
inertia conjugacy are exactly these Kummer classes; thus
each label supplies at most one full image.  An ordinary
pair \((x,1-x)\), \(x\ne0,1\), cannot equal any cusp
pair: the first two would give \(x=0\) or \(x=1\), and
the third would give \((1-x)/x=-1\).  All three equations
are excluded in the field by the same local injection.
Consequently \(\mathcal D_v\) has at most three elements,
or one when an ordinary image exists.

At the real place, ordinary points have at most the three
full section classes supplied by the three real intervals.
For precision, a path inside one interval transports every
finite etale fibre compatibly; its real transport commutes
with complex conjugation.  Taking the inverse limit of the
finite compatible transport maps gives the actual geometric
conjugacy of the full point sections.  Each cusp has only
the two possible real Kummer parameters, since
\(\widehat{\R^\times}=\{1,-1\}\).  Hence there are at
most nine candidate full images before identifying any
overlaps.  Formula (A.205.7) gives their cosets.  This
retains every real and cuspidal ambiguity required by the
actual local maps.

\emph{Fixed-support cohomology and compactness.}
Put \(m=r+1\) and \(M=\Z_\ell(m)\).  A continuous
one-cocycle on \(G_\Q\) restricts on inertia at
\(q\ne\ell\) to a homomorphism, since inertia acts
trivially on \(M\).  Wild and the other tame primary
factors map to zero.  Its value \(a\) on tame
\(\Z_\ell(1)\) satisfies
\(q^m a=q a\), by the cocycle and tame conjugation
identities.  Torsion-freeness gives \(a=0\).
The cocycle also vanishes on all conjugates of those
inertia groups, and on their closed normal closure.
Thus the actual inflation is an isomorphism
\[
 H^1_{\rm cts}(G_{\Q,\{\ell,\infty\}},M)
          \xrightarrow{\ \sim\ }H^1_{\rm cts}(G_\Q,M).
\tag{A.205.10}
\]
The identical local argument, followed by evaluation on
arithmetic Frobenius, gives
\[
 H^1(G_{\Q_q},\Z_\ell(m))
          =\Z_\ell/(q^m-1)\Z_\ell\quad(q\ne\ell).
\tag{A.205.11}
\]
Every Frobenius value occurs, by the compatible finite
geometric-sum cocycles on the procyclic unramified group,
as in A.60.  Thus these are the same actual local maps
used there, with the coefficient weight specified.

Finite fixed-support cohomology is finite by
\cite[I, Corollary 4.15]{BIR_MilneADT}.
The continuous comparison
\cite[(1.6), Theorem 2.2]{BIR_Jannsen1988} identifies
the group in (A.205.10) with the inverse limit of its
finite-coefficient \(H^1\) groups: the finite \(H^0\)
groups have zero derived inverse limit.  It is therefore
a compact Hausdorff pro-\(\ell\) group.
It is a finitely generated \(\Z_\ell\)-module: reduction
modulo \(\ell\) injects its quotient by \(\ell\) into
the finite \(H^1\) group with coefficients \(M/\ell\).
Choose lifts of generators.  Their closed span has
compact quotient \(V\) satisfying \(V=\ell V\).
Every finite pro-\(\ell\) quotient of \(V\) is then
zero, so \(V=0\).  This is the complete Nakayama argument
on the actual inverse-limit topology.

The same comparison gives compact Hausdorff local groups.
Taking the product over coefficient primes proves compactness
of \(\mathcal H\) and continuity of all its localization
maps.  Their simultaneous image is compact and hence closed
in the Hausdorff product of the local groups.  Each
\(\mathcal D_v\) is a finite union of the compact cosets
proved above, so its inverse image is closed in
\(\mathcal H\).  The finite-intersection property now
proves the asserted exact existence criterion.
Equivalently, the finite coefficient comparison maps
separate points of every local group; testing their finite
images at finitely many places and coefficient primes
defines the same closed intersection.  This justifies the
finite-power formulation without replacing any image by
its closure.

\emph{The exact fourth-weight global group.}
Take \(M=\Z_\ell(4)\) and \(G_S=G_{\Q,\{\ell,\infty\}}\).
A.61 proves \(H^2(G_S,M)=0\) at odd \(\ell\) and an
injection into \(\Z/2\) at two.  Moreover
\(H^3_{\rm cts}(G_S,M)=0\).
Indeed finite coefficient \(H^3\) is real cohomology by
\cite[I, Theorem 4.10(c)]{BIR_MilneADT}; at two the even
twist gives \((M/2^a)[2]\), and the reduction maps from
level \(a+1\) to \(a\) are zero.  The inverse limit is
zero, and the finite \(H^2\) groups have zero derived
inverse limit.  At odd coefficients the real group is zero.

The coefficient sequence therefore bounds
\(H^2(G_S,M/\ell^a)\) by order two at \(\ell=2\), and
makes it zero otherwise.  The exact Euler-characteristic
formula \cite[I, Theorem 5.1]{BIR_MilneADT}, with its
ordinary real \(H^0\) term and the even twist, gives
\[
 \#H^1(G_S,M/\ell^a)=
       \#H^0(G_S,M/\ell^a)\,\#H^2(G_S,M/\ell^a).
\tag{A.205.12}
\]
This is the same real normalization as (A.40.10), not
an assumption of vanishing higher real cohomology.
The cyclotomic image on \(G_S\) is \(\Z_\ell^\times\),
and its exact fourth-power ideals are
\[
 (u^4-1:u\in\Z_\ell^\times)=
 \begin{cases}
 16\Z_2,&\ell=2,\\
 3\Z_3,&\ell=3,\\
 5\Z_5,&\ell=5,\\
 \Z_\ell,&\ell\geq7.
 \end{cases}
\tag{A.205.13}
\]
Odd fourth powers are one modulo sixteen and \(3^4-1=80\)
attains that two-adic valuation.  At three and five,
\(2^4-1=15\) attains valuation one.  For \(\ell\geq7\),
the degree-four polynomial has at most four roots among
the at least six nonzero residue classes, giving a unit.
It follows that the finite \(H^1\) groups have orders
bounded by \(32,3,5,1\), respectively.  Their inverse
limits are finite.

For any finite integral \(H^1\) class \([c]\), choose
\(a\) with \(\ell^a[c]=0\).  Then
\(\ell^ac(g)=(\chi_\ell(g)^4-1)t\) for some
\(t\in\Z_\ell\).  Division inside \(\Q_\ell(4)\)
shows that \(c\) is the fractional coboundary of
\(t/\ell^a\).  Conversely an invariant fractional class
gives an integral cocycle.  The absence of invariants in
\(\Q_\ell(4)\) proves the exact isomorphism
\[
 H^0(G_S,(\Q_\ell/\Z_\ell)(4))
   \xrightarrow{\ \sim\ }H^1_{\rm cts}(G_S,\Z_\ell(4)),
 \quad \bar t\longmapsto[(\chi_\ell^4-1)t].
\tag{A.205.14}
\]
Its orders from (A.205.13) are \(16,3,5,1\).
Taking the actual product of the coefficient cochain
complexes proves (A.205.3), with its stated generator.
There are three independent fourth-weight Hall coordinates,
so the complete lift torsor has \(240^3\) members.

Restriction of a nonzero fractional-coboundary class to
any finite place remains nonzero.  Indeed, if the restriction
of \(k\eta_4\) were the integral coboundary of \(t\),
then for every local \(g\)
\[
       (\chi_\ell(g)^4-1)(k/240-t)=0.
\tag{A.205.15}
\]
Some such factor is nonzero at every coefficient prime,
by the same local cyclotomic argument as above.  Thus
\(k/240=t\in\Z_\ell\); for all coefficients this forces
\(k=0\pmod {240}\).  This proves restriction injectivity.
For \(q\ne\ell\), its actual Frobenius value is
\((q^4-1)k/240\) modulo \((q^4-1)\Z_\ell\).
Putting \((q,\ell)=(3,2),(2,3),(2,5)\) gives precisely
(A.205.4), which is an isomorphism by the Chinese
remainder theorem.  Applying it in each Hall coordinate
proves the asserted partial normalization of any global
lift supplied by A.204.

Finally \(H^1(G_\R,\widehat\Z(4)^3)=0\), since the
even twist is trivial and the coefficient module has no
two-torsion.  Every global class-four lift of a genuinely
real-matching lower section therefore has the same local
full point or cusp image at this level.  In this case
(A.205.7) has zero target and no hidden real orbit remains.

If every candidate in the finite global torsor fails
some local condition, choose one failing place per
candidate.  The union has at most \(240^3\) members
and already makes the matching intersection empty.
At a finite place with an ordinary lower image,
\(\mathcal D_v\) is a singleton and restriction is
injective, so there is at most one surviving candidate.
If it exists but does not match everywhere, one more
failing place excludes it.  This proves the sharper
two-place assertion and completes the exact local
correction criterion.
\end{proof}
